
\section{Métodos propuestos}\label{sec:mp}
\begin{spacing}{1.5}

La presente investigación adopta un enfoque cuantitativo y experimental para el desarrollo de un generador paramétrico. El proceso se divide en la selección de recursos, el diseño del modelo matemático y el algoritmo de ajuste.

\subsection{Recursos Computacionales y Software}

Para garantizar la reproducibilidad y eficiencia del generador se utilizaron los siguientes recursos:

\begin{itemize}
    \item \textbf{Lenguaje de Programación:} Python 3.12, debido a su robustez en el procesamiento de señales y optimización numérica.

    \item \textbf{Librerías Núcleo:} NumPy y SciPy \cite{virtanen2020scipy}: Para el manejo de arreglos vectoriales y la implementación de algoritmos de optimización.

    \begin{itemize}
        \item WFDB (Waveform Database): Para la lectura y preprocesamiento de los registros de la base de datos MIT-BIH Arrhythmia Database \cite{moody2001impact}.
        \item Matplotlib / Seaborn: Para la visualización y comparación de señales sintéticas frente a reales.
    \end{itemize}

    \item \textbf{Hardware}: Estación de trabajo con procesador multi-núcleo y un mínimo de 8 GB de RAM, suficiente para la ejecución de modelos paramétricos gaussianos.
\end{itemize}

Los resultados reportados se obtuvieron con Python 3.12.4, NumPy 2.1.3 y SciPy 1.16.0. Registrar las versiones no es una formalidad. El ajuste morfológico es determinista: no interviene ningún muestreo aleatorio y el optimizador parte de un punto fijo por latido, de modo que repetir el ajuste en el mismo entorno reproduce exactamente todos los parámetros. El módulo de ritmo sí emplea números aleatorios, y recibe la semilla como parámetro explícito. Sin embargo, al reajustar los mismos latidos con versiones distintas de las bibliotecas de cálculo, L-BFGS-B converge a puntos ligeramente diferentes: sobre los primeros cuarenta latidos del registro, el error cuadrático medio difirió en torno al 1\,\%, en una dirección u otra según el latido. El procedimiento es, por lo tanto, determinista dentro de un entorno y no entre entornos. Por esa razón cada archivo de métricas que produce el pipeline registra el entorno con que se generó.

\subsection{Diseño del Modelo de Generación}

El núcleo del generador se basa en el Modelado Fenomenológico Paramétrico, que describe la señal ECG como una combinación lineal de funciones matemáticas.

\subsubsection{Modelo de Gaussianas Asimétricas}

Cada ciclo cardíaco (latido) se modela mediante la superposición de cinco funciones gaussianas asimétricas, correspondientes a las ondas P, Q, R, S y T, más un término de línea base afín. La señal sintética $x(t)$ se define como:

\begin{equation}
    x(t) = \sum_{i \in \{P,Q,R,S,T\}} a_i \exp\!\left( -\frac{(t - b_i)^2}{2\,c_i(t)^2} \right) \;+\; \beta_0 + \beta_1\,(t - t_R)
\end{equation}
donde el ancho de cada onda depende del lado de su centro en que se evalúe:

\begin{equation}
    c_i(t) =
    \begin{cases}
        c_i^{\text{asc}} & \text{si } t < b_i \\[2pt]
        c_i^{\text{desc}} & \text{si } t \geq b_i
    \end{cases}
\end{equation}

Cada onda queda descrita por cuatro parámetros explícitos: su amplitud ($a_i$), la posición de su centro ($b_i$) y los anchos de sus ramas ascendente ($c_i^{\text{asc}}$) y descendente ($c_i^{\text{desc}}$). El término de línea base aporta dos más, el desnivel $\beta_0$ y la pendiente $\beta_1$, referidos al instante $t_R$ del pico R. El total es de veintidós parámetros por latido.

La asimetría no es un refinamiento estético. La onda T sube más rápido de lo que baja, y una gaussiana simétrica no puede representar esa forma por sí sola: hacerlo exigiría superponer una segunda onda, con tres parámetros más sin significado fisiológico propio. Con dos anchos por onda, la asimetría se describe con un único parámetro adicional que conserva una interpretación directa.

\subsection{Proceso de Optimización Numérica}

Para que las señales sean realistas, los parámetros del modelo se ajustan automáticamente a señales reales utilizando técnicas de optimización. El flujo de trabajo consiste en:

\begin{enumerate}
    \item \textbf{Extracción de Semillas:} Selección de latidos representativos de la base MIT-BIH.
    \item \textbf{Función Objetivo:} Minimización del error cuadrático medio (MSE) entre la señal real ($y$) y la sintética generada ($x$). Se minimiza el MSE y se reporta su raíz, el RMSE, por las razones que se detallan en la Subsección \ref{sec:metricas}.
    \item \textbf{Algoritmo:} Se utiliza el método \textbf{L-BFGS-B} \cite{byrd1995limited}, permitiendo establecer límites fisiológicos estrictos a los parámetros para evitar artefactos biológicamente imposibles.
\end{enumerate}

\subsection{Motor de Reglas Fisiológicas y Variabilidad}

El diseño contempla, para la etapa siguiente del plan de trabajo, un motor de reglas basado en estándares clínicos que permita generar dos grupos etiquetados:

\begin{itemize}
    \item \textbf{Grupo control (sanos):} Ajuste automático de los intervalos (PR, QRS, QT) y de las amplitudes dentro de los rangos normales de referencia establecidos por las tablas de Rijnbeek \cite{rijnbeek2001normal} y la AHA \cite{rautaharju2009aha}, con una serie de intervalos R-R cuyo exponente de escalamiento corresponde al régimen reportado para sujetos sanos.
    \item \textbf{Grupo de enfermos:} Desvío controlado de esos mismos parámetros fuera de los rangos normales de referencia, acompañado de una serie de intervalos R-R con exponente de escalamiento degradado. La etiqueta es de construcción y no diagnóstica: registra cómo se generó el sujeto sintético, no una condición clínica atribuida.
\end{itemize}


\subsection{Metodología de validación de los datos generados}

Esta sección responde a la observación de la comisión relativa a la necesidad de definir explícitamente la metodología de aplicación de las métricas de validación. Se especifican la unidad de análisis, el procedimiento de comparación, las métricas empleadas con su justificación, y los criterios de aceptación.

\subsubsection{Unidad de análisis y procedimiento de comparación}

La unidad de validación es el latido individual. Cada señal sintética se compara exclusivamente con el latido real que sirvió de semilla para su ajuste, de modo que la comparación es pareada y no agregada. Se descarta deliberadamente la comparación contra un latido promedio del registro, dado que el promediado atenúa la variabilidad morfológica que el generador debe reproducir.

El procedimiento consta de cuatro etapas:

\begin{enumerate}
    \item \textbf{Selección de semillas.} Extracción de los latidos anotados como sinusales normales (símbolo \texttt{N}) del registro, descartando latidos ectópicos, artefactos y marcas de ritmo. Del registro 100 se obtienen así 2.237 latidos de las 2.239 anotaciones normales. Los dos restantes se descartan por motivos distintos: el último latido del registro, porque el recorrido de las anotaciones excluye la primera y la última como precaución frente a los bordes —y en este registro la primera es un marcador de ritmo, no un latido—, y un latido cuya ventana de extracción comienza antes del inicio de la señal.
    \item \textbf{Segmentación y normalización.} Ventana asimétrica respecto del pico R anotado, de 250 ms antes y 500 ms después, es decir 270 muestras a la frecuencia de muestreo de 360 Hz del registro. La asimetría se sigue de la fisiología del complejo P-QRS-T, y el efecto de haber reemplazado la ventana centrada que se empleó en una iteración anterior se cuantifica en la Subsección \ref{sec:validez}. Cada latido se normaliza mediante puntuación Z, de modo que las métricas resulten independientes de la ganancia del canal y comparables entre registros.
    \item \textbf{Ajuste.} Optimización independiente de los veintidós parámetros del modelo para cada latido, minimizando el error cuadrático medio mediante L-BFGS-B \cite{byrd1995limited} con restricciones de caja fisiológicas.
    \item \textbf{Evaluación.} Cálculo de las métricas sobre el par formado por el latido real y su reconstrucción sintética.
\end{enumerate}

\subsubsection{Métricas empleadas y justificación de su elección}\label{sec:metricas}

Se emplean dos métricas complementarias, cuya combinación responde a que ninguna resulta suficiente por separado.

La \textbf{correlación de Pearson} evalúa la concordancia morfológica entre ambas señales:

\begin{equation}
r = \frac{\sum_{i=1}^{N}(y_i - \bar{y})(x_i - \bar{x})}{\sqrt{\sum_{i=1}^{N}(y_i - \bar{y})^2}\sqrt{\sum_{i=1}^{N}(x_i - \bar{x})^2}}
\end{equation}

donde $y$ denota la señal real, $x$ la sintética y $N = 270$ el número de muestras de la ventana. Se selecciona porque el objetivo primario del generador es reproducir la forma del ciclo P-QRS-T. Su limitación conocida es la invarianza ante transformaciones afines, motivo por el cual no se reporta de forma aislada.

El \textbf{error cuadrático medio y su raíz} cuantifican la discrepancia punto a punto:

\begin{equation}
\mathrm{MSE} = \frac{1}{N}\sum_{i=1}^{N}(y_i - x_i)^2 \qquad ; \qquad \mathrm{RMSE} = \sqrt{\mathrm{MSE}}
\end{equation}

Se distingue explícitamente entre ambas magnitudes: el MSE constituye la función objetivo minimizada durante la optimización, mientras que el RMSE es la métrica reportada, por estar expresada en las mismas unidades que la señal normalizada y resultar por ello directamente interpretable. Esta distinción se hace explícita porque ambas denominaciones aparecen empleadas de forma intercambiable en la literatura del área, lo que constituye una fuente frecuente de confusión.

\subsubsection{Estadísticos reportados}

Para cada métrica se reportan media, desviación estándar, mínimo, mediana y máximo sobre el total de latidos procesados. La desviación estándar y el rango son tan relevantes como la media: un valor promedio elevado acompañado de dispersión alta indicaría que el generador funciona bien únicamente sobre un subconjunto de latidos, situación que la media por sí sola ocultaría.

Se reporta asimismo la proporción de latidos que supera umbrales de correlación definidos, de modo que el desempeño pueda evaluarse en términos de cobertura y no solo de tendencia central.

\subsubsection{Criterios de aceptación}

Se establecen los siguientes criterios, definidos con anterioridad a la ejecución de los experimentos:

\begin{itemize}
    \item Correlación de Pearson media superior a 0,95, con desviación estándar inferior a 0,05.
    \item Al menos el 90\% de los latidos con correlación individual superior a 0,95.
    \item Ausencia de latidos con parámetros fuera de los rangos fisiológicos definidos. Un parámetro que converge exactamente sobre su cota se considera dentro del rango, pero no se da por determinado por los datos: esa situación se cuantifica aparte, mediante el control de saturación que se describe a continuación.
\end{itemize}

\subsubsection{Verificación de saturación de las restricciones}

Como control adicional del procedimiento de optimización se verifica, para cada parámetro, la proporción de latidos cuya solución converge exactamente sobre alguno de los límites de caja impuestos. Una proporción elevada indica que la restricción, y no los datos, está determinando el valor del parámetro, situación que invalida la interpretación fisiológica de dicho parámetro y exige recalibrar el límite correspondiente. Este control forma parte del protocolo de validación y sus resultados se reportan en la sección correspondiente.

\subsubsection{Alcance de la validación}

La validación descrita se aplica a señales de ritmo sinusal normal. La evaluación de señales con variabilidad demográfica parametrizada o con alteraciones patológicas requerirá criterios adicionales, específicamente el contraste de los intervalos generados contra los rangos de referencia clínicos de Rijnbeek \cite{rijnbeek2001normal} y de la AHA \cite{rautaharju2009aha}, y queda fuera del alcance de la presente etapa.

\subsection{Módulo de ritmo y acoplamiento con la morfología}

El modelo descrito hasta aquí produce latidos. Esta sección describe cómo se
produce la secuencia temporal en que esos latidos se emiten, y cómo ambas
partes se unen.

\subsubsection{Generación de la serie de intervalos}

La serie de intervalos R-R se sintetiza en el dominio de la frecuencia. Se
construyen amplitudes proporcionales a $f^{-\beta/2}$, se les asignan fases
aleatorias uniformes y se antitransforma, de modo que la serie resultante tenga
densidad espectral de potencia proporcional a $f^{-\beta}$. Dado que
$\alpha = (\beta+1)/2$, pedir un exponente de escalamiento $\alpha$ fija el
exponente espectral $\beta$, y el procedimiento cubre de manera continua el
rango de $\alpha$ entre 0,5 y 1,5.

Dos detalles del procedimiento son deliberados. Primero, la transformada
inversa impone periodicidad circular, lo que introduce una correlación espuria
entre el comienzo y el final de la serie; la mitigación estándar consiste en
sintetizar una serie varias veces más larga que la requerida y recortarla, y
aquí se emplea un factor de cuatro. Segundo, la serie normalizada se escala a la
desviación estándar pedida y se le suma el intervalo medio, y el resultado se
acota al rango fisiológico. Ese recorte se reporta: si algún intervalo debió
recortarse, la combinación de dispersión y exponente solicitada no es
fisiológicamente sostenible y el exponente efectivo de la serie ya no es el
pedido. Recortar en silencio equivaldría a falsear el resultado.

Se parametriza en $\alpha$ y no en el exponente de Hurst $H$ por una razón
empírica: los valores publicados para sujetos sanos en registros de
veinticuatro horas se ubican en torno a 1,17
\cite{costa2017fragmentation}, es decir por encima de la unidad, y por lo tanto
fuera del rango que un generador de ruido gaussiano fraccional puede producir.

\subsubsection{Medición del exponente}

El exponente de la serie generada se mide con una implementación propia del DFA,
según el procedimiento de Peng et al.\ \cite{peng1995quantification}. El
exponente global se estima sobre ventanas de dieciséis muestras en adelante,
porque por debajo de esa escala el estimador sobreestima el exponente incluso
sobre series construidas sin memoria.

Se reportan además los exponentes por tramos de uso clínico, $\alpha_1$ sobre
ventanas de 4 a 16 latidos y $\alpha_2$ sobre ventanas de 16 a 64, dado que el
cruce de pendiente descrito en la literatura hace que un único exponente global
pueda ocultar comportamientos distintos en escalas cortas y largas. El primero
se reporta junto a su sesgo, porque se calcula precisamente en el rango de
escalas donde el estimador sobreestima. Medido con esta misma implementación
sobre ruido gaussiano de exponente verdadero 0,5, con sesenta repeticiones por
cada largo de serie, $\alpha_1$ arroja valores entre 0,587 y 0,591 para series
de 512 a 32.768 puntos. El sesgo es de aproximadamente $+0{,}09$ y no disminuye
con el largo del registro, de modo que es atribuible al estimador en ventanas de
pocos puntos y no a la serie. Se conserva el rango de 4 a 16 por ser la
convención clínica, pero ningún valor de $\alpha_1$ se informa sin esa
advertencia.

\subsubsection{Acoplamiento selectivo en el tiempo}

Unir morfología y ritmo exige decidir cómo cambia la forma del latido cuando
cambia su duración. Mantener la morfología fija sería incoherente: un latido a
120 latidos por minuto tendría la misma duración de onda T que uno a 50. Y
estirar el latido completo en proporción al intervalo también lo sería, porque
la duración del complejo QRS depende de la velocidad de conducción
intraventricular y no de la frecuencia cardíaca.

El modelo paramétrico permite una tercera opción, que es la adoptada: escalar el
tiempo de forma \emph{selectiva}. Los parámetros de las ondas Q, R y S se dejan
intactos, mientras que los centros y los anchos de las ondas P y T se escalan
por $(RR/RR_{\text{ref}})^{\gamma}$. Esta posibilidad es una ventaja estructural
del enfoque paramétrico frente a los modelos de ecuaciones diferenciales, donde
la velocidad angular gobierna el ciclo completo y todas las ondas se estiran en
bloque.

El valor por defecto es $\gamma = 1$, la relación proporcional, por dos motivos.
Es la única relación que se pudo verificar en la bibliografía disponible:
McSharry et al.\ \cite{mcsharry2003dynamical}, citando a Davey
\cite{davey1999qt}, afirman que el intervalo QT crece linealmente con el
intervalo R-R. Y reproduce el comportamiento implícito del generador de
referencia, lo que permite una comparación directa. Se declara, no obstante, una
aproximación: una relación lineal con ordenada no nula es más plana que una
proporcional, de modo que $\gamma = 1$ sobrecorrige el intervalo QT a
frecuencias bajas. El exponente se expone como parámetro para poder contrastar
las correcciones alternativas de uso clínico.

\subsubsection{Construcción de la señal continua}

Cada latido se evalúa sobre su propio tramo de tiempo, delimitado por los puntos
medios respecto de los picos R vecinos, y los tramos se unen con un cruce suave
de veinte milisegundos para que la diferencia de nivel entre latidos vecinos no
deje un escalón que el análisis espectral leería como contenido de alta
frecuencia.

La línea base recibe un tratamiento aparte. Sus dos parámetros se ajustaron
dentro de la ventana de un latido, y extrapolar esa tendencia afín a lo largo de
un intervalo R-R completo multiplica el desnivel y produce un diente de sierra.
En su lugar, el nivel de línea base ajustado de cada latido en su pico R se toma
como nudo y se interpola entre latidos con un polinomio cúbico monótono por
tramos, lo que entrega la deriva lenta y continua que caracteriza a una línea
base real.

\subsubsection{Criterio de validación del módulo}

El estado del arte en generación de series de intervalos validó por
indistinguibilidad: en el desafío de 2002 \cite{moody2002challenge} los
participantes debían separar series reales de sintéticas en un conjunto sin
etiquetar. Este trabajo adopta un criterio distinto y verificable de manera
independiente: se pide al generador un exponente, se mide el exponente de la
serie que produjo, y se caracteriza el sesgo del estimador en función del largo
del registro. La validación no consiste en que la serie parezca real, sino en que
posea, de manera medible, la propiedad que se le pidió.

\end{spacing}
